\chapter{环境与健康的关系}
\setcounter{chapter}{2}
\chapterintro{
本章是环境卫生学的理论基础，共12学时（占比最大）。掌握生态系统、食物链、生物放大等基本概念；掌握人群健康效应谱、剂量-反应关系；掌握环境流行病学和毒理学研究方法的特点；掌握健康危险度评价的四个阶段。
}

% ----- 2.1 人类的环境 -----
\section{人类的环境}

\subsection{生态系统}
\begin{defbox}
\textbf{生物圈（Biosphere）}：从海平面以下深约12 km至海平面以上高约10 km的范围，包括一部分大气圈、水圈和土壤岩石圈，是地球上所有生命物质及其生存环境的整体。
\end{defbox}

\begin{defbox}
\textbf{生态系统（Ecosystem）}：在一定空间范围内，由\keyword{生物群落及其环境}组成，借助各种功能流（\imp{物质流、能量流、物种流和信息流}）所联结的稳态系统。
\end{defbox}

\textbf{生态系统组成：}
\begin{itemize}[leftmargin=*]
    \item \imp{生产者}（绿色植物、光合细菌）——将无机物转化为有机物
    \item \imp{消费者}（动物）——直接或间接以生产者为食
    \item \imp{分解者}（微生物）——将有机物分解为无机物
    \item \imp{无机界}（阳光、水、空气、土壤等）
\end{itemize}

\textbf{生态系统的主要特征：}
\begin{enumerate}[leftmargin=*]
    \item 整体性——各要素对系统整体性起作用
    \item 开放性——与外界进行物质交换
    \item 自调控——种群密度调控、种间数量调控、生物与环境适应调控
    \item 可持续性——不断进行物质循环和能量流动
\end{enumerate}

\begin{defbox}
\textbf{生态系统健康（Ecosystem Health）}：具有活力、结构稳定和自调节能力的生态系统，是生态系统的综合特性。
\end{defbox}

\subsection{食物链与生物放大}

\begin{defbox}
\textbf{食物链（Food Chain）}：生态系统中不同营养级的生物逐级被吞食以满足生存需要而建立起来的链锁关系。
\end{defbox}

\begin{defbox}
\textbf{生物富集作用（Bioenrichment）}：生物从环境中摄入浓度极低的重金属元素或难降解化合物，在体内逐渐累积，使生物体内该物质的浓度大大超过环境中的浓度。

\textbf{浓集系数} = 污染物在生物体内浓度 / 污染物在环境中浓度。
\end{defbox}

\begin{defbox}
\textbf{迁移（Transportation）}：环境物质在环境中发生的空间位移过程，伴随污染物在环境介质中浓度改变。

\textbf{转化（Transform）}：化学物（或污染物）在环境中通过物理、化学或生物作用改变其存在形态或转变成另一物质的过程。
\end{defbox}

\begin{keybox}
\textbf{环境污染物的迁移和转化对环境因素暴露的影响：}
\begin{enumerate}[leftmargin=*]
    \item 扩大暴露范围
    \item 增加暴露途径
    \item 影响暴露剂量
    \item 改变污染物性质和毒性
\end{enumerate}

\textbf{生物放大作用（Biomagnification）}：环境中的重金属元素和难降解有机物，通过食物链转移到高位营养级生物体内，使其浓度逐级在生物体内放大。

\textbf{典型实例：}DDT在水生食物链中的迁移——浮游生物（0.04 ppm）→ 小鱼（0.5 ppm）→ 大鱼（2.0 ppm）→ 水鸟（25 ppm），浓缩系数达600倍以上。

\textbf{水俣病}和\textbf{痛痛病}就是生物放大作用的经典案例。
\end{keybox}

% ----- 2.2 人与环境的辩证关系 -----
\section{人与环境的辩证统一关系}

\subsection{人与环境在物质上的统一性}
人体血液中的60多种元素含量与地壳中这些元素的含量呈明显正相关，说明人体组成与地球化学环境具有\keyword{物质上的统一性}。

\subsection{人类对环境的适应性}
\begin{itemize}[leftmargin=*]
    \item 光适应
    \item 高原低氧适应
    \item 热适应
    \item 人类进化适应
\end{itemize}

\textbf{生物学基础：}环境应答基因（Environmental Response Gene）的多态性是造成人群易感性差异的重要原因。涉及\imp{基因多态性}和\imp{单核苷酸多态性（SNP）}。

\begin{defbox}
\textbf{环境基因交互作用}：指个体基因与外部环境相互作用的过程。①环境影响基因的表达——如某些环境化学物既是CYP酶的底物又是酶的诱导剂，可增强CYP酶的代谢能力，从而增毒或解毒；②基因多态性影响有害因素的作用方式——如影响化学物在体内的吸收、转运、分布、代谢等。

典型实例：\textbf{苯丙酮尿症}——苯丙氨酸羟化酶缺乏且摄入苯丙氨酸；\textbf{肿瘤发生}——致癌因素暴露 + DNA修复基因SNP + 代谢酶遗传多态性。

机体适应分为\imp{群体适应}（通过基因或表观遗传改变实现，可代际遗传）和\imp{个体适应}（不涉及遗传物质改变，如个体对高温、高原的适应）。
\end{defbox}

\subsection{环境因素对健康影响的双重性}
环境因素既有\imp{有益作用}（如紫外线促进维生素D合成），也有\imp{有害作用}（如过量紫外线导致皮肤癌）。

% ----- 2.3 暴露特征 -----
\section{环境因素的暴露特征}

\subsection{暴露的基本概念}

\begin{defbox}
\textbf{暴露（Exposure）}：人体接触某一有害环境因素的过程。
\end{defbox}

\begin{itemize}[leftmargin=*]
    \item \imp{外剂量（External Dose）}：环境介质中某种环境因素的浓度或含量，根据人体接触特征估计个体的暴露水平。
    \item \imp{内剂量（Internal Dose）}：机体内已吸收的污染物的量。
    \item \imp{生物有效剂量（Biologically Effective Dose）}：经吸收、代谢活化、转运最终到达靶部位或替代性靶部位的污染物量。
\end{itemize}

\begin{defbox}
\textbf{生物半减期（Biological Half-life, t$_{1/2}$）}：化合物在体内含量减少一半所需的时间。化合物进入机体经历六个生物半减期后，体内最大可能蓄积量趋于稳定，此后摄入量与排出量趋于平衡。摄入量愈大，达到平衡后其最大蓄积量也愈大；摄入量若减少到致使体内最大蓄积量低于产生有害效应的水平，长期作用也观察不到对机体的有害效应。
\end{defbox}

\subsection{环境多因素暴露与联合作用}

\begin{table}[htbp]
\centering
\caption{联合作用的类型}
\begin{tabularx}{\textwidth}{lX}
\hline
\textbf{类型} & \textbf{定义及举例} \\
\hline
\imp{相加作用（Additive）} & 多种化学物联合作用强度等于各化学物单独作用强度之和。如同系物或毒作用靶器官相同的化学物。大部分刺激性气体的刺激作用、两种有机磷农药抑制胆碱酯酶的作用 \\
\imp{独立作用（Independent）} & 各化学物作用于不同受体/部位/靶器官，引发的生物效应互不干扰，表现为各自毒性效应。如酒精与氯乙烯的联合作用 \\
\imp{协同作用（Synergistic）} & 联合作用强度大于各化学物单独作用强度之和。如马拉硫磷和苯硫磷的联合作用 \\
\imp{增强作用（Potentiation）} & 一种化学物本身对某器官无毒性，但能增强另一种化学物的毒性。如异丙醇对肝脏无毒性，但与四氯化碳同时进入机体时使四氯化碳毒性远高于单独作用 \\
\imp{拮抗作用（Antagonism）} & 联合作用强度小于各化学物单独作用强度之和。包括：
\begin{itemize}
    \item 竞争作用：肟类化合物与有机磷竞争胆碱酯酶
    \item 功能性拮抗：阿托品对抗有机磷的毒蕈碱症状
    \item 代谢过程变化：卤代苯类诱导有机磷代谢使毒性减弱
\end{itemize} \\
\hline
\end{tabularx}
\end{table}

\subsection{剂量-反应关系}

\begin{defbox}
\textbf{无阈值化合物（No-threshold Compound）}：在大于零的剂量暴露下均可能发生有害效应的化合物。其剂量-反应曲线延长线通过坐标原点，\textbf{无安全剂量}。典型代表：遗传毒性致癌物。

\textbf{阈值化合物（Threshold Compound）}：仅在达到或大于阈剂量时才产生效应。剂量-反应曲线多呈\textbf{S形或抛物线形}。有两个阈值的化合物呈\textbf{U形}曲线。
\end{defbox}

\begin{keybox}
\textbf{阈值理论在环境卫生标准制定中的重要性：}对于有阈化合物，可通过毒理学和流行病学方法确定阈剂量，进而除以安全系数得到卫生标准限值。无阈化合物（如遗传毒性致癌物）则不适用阈值理论，需采用健康风险评价方法。
\end{keybox}

\begin{tipbox}
\textbf{研究环境因素联合作用的意义：}
\begin{enumerate}[leftmargin=*]
    \item 多种环境因素同时存在时对人体的作用与其中一种单独存在时有所不同，往往呈现复杂的交互作用，使机体的毒性效应发生改变
    \item 在阐明环境因素对健康的影响、制订混合污染物的卫生标准、开展多种有害因素共同作用的危险度评价、采取防治对策等方面具有实际意义
\end{enumerate}
\end{tipbox}

% ----- 2.4 健康效应 -----
\section{环境因素的健康效应}

\subsection{人群健康效应谱}

\begin{keybox}
\textbf{人群健康效应谱（Spectrum of Health Effect）}：人群暴露于环境污染物时产生的不良反应在性质、程度和范围上表现的征象，从弱到强分为五级：

\begin{enumerate}[leftmargin=*]
    \item \textbf{污染物负荷增加}——体内负荷增加，不引起生理功能和生化代谢变化
    \item \textbf{生理代偿性变化}——出现某些生理功能和生化代谢变化（生理代偿）
    \item \textbf{亚临床状态（准病态）}——生理代谢异常，机体处于病理性代偿调节状态，无明显临床症状
    \item \textbf{临床性疾病}——机体功能失调，出现临床症状
    \item \textbf{死亡}——严重中毒导致死亡
\end{enumerate}

\textbf{分布规律：}呈金字塔形分布——健康人群最多，亚临床状态居中，临床患者较少，死亡最少。
\end{keybox}

\begin{defbox}
\textbf{高危人群（High Risk Group）}：对环境暴露的健康效应出现较早、效应较显著、易感性较高的人群，包括高暴露人群（职业人群和特殊暴露人群）和高敏感人群（易感人群）。

\textbf{影响人群易感性的因素：}
\begin{itemize}[leftmargin=*]
    \item \imp{非遗传因素}：年龄、健康状况、营养状态、生活习惯、暴露史、心理状态、保护措施
    \item \imp{遗传因素}：性别、种族、遗传缺陷、基因多态性
\end{itemize}

\textbf{脆弱人群（Vulnerable Group）}：身体状况较差或抗病能力较弱的人。

\textbf{敏感人群（Sensitive Group）}：接触有害物质时，由于个体生物学因素使其毒性反应的出现较普通人群更早、反应更强的人群。
\end{defbox}

\begin{tipbox}
\textbf{研究高危人群的意义：}
\begin{enumerate}[leftmargin=*]
    \item 是最早出现健康效应的人群，有助于全人群健康效应的早发现和早预防
    \item 高危人群的健康效应更复杂和严重，对该人群的保护即是对全人群的保护
    \item 对高危人群提早筛查可获得最大流行病学和卫生经济学效益
\end{enumerate}
\end{tipbox}

\subsection{环境污染对人群健康的危害}

\begin{enumerate}[leftmargin=*]
    \item \imp{急性危害（Acute Hazard）}：污染物短时间内大量进入环境，使暴露人群短时间内出现不良反应、急性中毒甚至死亡。包括：大气污染烟雾事件、过量排放和事故性排放、生物性污染引起的急性传染病。
    \item \imp{慢性危害（Chronic Hazard）}：有害污染物以低浓度、长时间反复作用于机体产生的危害。包括非特异性影响、慢性疾患（COPD等）、持续性蓄积危害。
    \item \imp{致癌致畸危害}：如反应停（Thalidomide）事件——孕妇服用导致胎儿畸形。
    \item \imp{环境致畸物和致畸因素}：孕期摄入环境污染物可引起胎儿畸形发生率增加。Wilson认为导致出生缺陷率增加的因素：遗传因素占25\%，环境因素占10\%，遗传与环境相互作用及原因不明者占65\%。已知人类致畸因素：辐射（原子武器、放射性碘）；感染（巨细胞病毒、水痘病毒、梅毒螺旋体）；母体代谢失衡（酒精中毒、糖尿病、叶酸缺乏）；药物和环境化学物（氨蝶呤、环磷酰胺、二噁英、氯乙烯等）。
\end{enumerate}

\subsection{致癌物分类}

\begin{table}[htbp]
\centering
\caption{IARC致癌物分类}
\begin{tabularx}{\textwidth}{llX}
\hline
\textbf{类别} & \textbf{分级} & \textbf{说明} \\
\hline
1类 & 对人致癌 & 流行病学证据充分（严格设计、排除混杂、有剂量-反应关系、有独立验证）\\
2A类 & 对人很可能致癌 & 动物实验证据充分，对人体致癌证据有限 \\
2B类 & 对人可能致癌 & 动物实验证据不充分或人体证据不足但动物证据充分 \\
3类 & 尚无法分类 & 可疑对人致癌 \\
4类 & 对人很可能不致癌 & — \\
\hline
\end{tabularx}
\end{table}

\begin{defbox}
\textbf{持久性有机污染物（POPs）}：具有持久性、蓄积性、迁移性和高毒性的有机污染物，可产生急慢性毒性及"三致"效应。

\textbf{环境内分泌干扰物（EDCs）}：对维持机体内环境稳态和调节发育过程的体内天然激素生成、释放、转运、代谢、结合、效应造成严重影响的一类外源性物质。与生殖障碍、出生缺陷、发育异常、代谢紊乱及某些癌症（乳腺癌、睾丸癌、卵巢癌等）的发生发展有关。
\end{defbox}

% ----- 2.5 研究方法 -----
\section{环境与健康关系的研究方法}

\subsection{环境流行病学方法}

\begin{keybox}
\textbf{环境流行病学（Environmental Epidemiology）}：应用传统流行病学方法，结合环境与人群健康关系的特点，从宏观上研究外环境因素与人群健康关系的科学。

\textbf{研究特点：}\imp{多因素、弱效应、非特异性、潜隐期}。

\textbf{基本内容：}
\begin{enumerate}[leftmargin=*]
    \item 研究已知环境因素对人群的健康效应
    \item 探索引起健康异常的环境有害因素
    \item 暴露剂量-反应关系的研究
\end{enumerate}
\end{keybox}

\textbf{（一）暴露测量}

\begin{itemize}[leftmargin=*]
    \item \imp{外暴露剂量}：通过测定人群接触的环境介质中某因素的浓度或含量，根据人体接触特征估计个体暴露水平。减少误差的方法：个体空气采样器、扩大样本量、综合考虑多种暴露途径估计总暴露量。
    \item \imp{内剂量（Internal Dose）}：已吸收至人体内的污染物量。通过测定生物材料（血液、尿液等）中污染物或其代谢产物的含量确定。
    \item \imp{生物有效剂量（Biologically Effective Dose）}：经吸收、代谢活化、转运最终到达靶部位或替代性靶部位的污染物量。
\end{itemize}

\textbf{（二）健康效应测量}

\begin{itemize}[leftmargin=*]
    \item 疾病频率测量：发病率、患病率、死亡率，各种疾病专率
    \item 生理生化功能测量：按手段分（生理、生化、分子生物学、遗传学、影像学）；按器官系统分（呼吸、消化、神经系统等）
\end{itemize}

\textbf{（三）因果关系判断标准}

关联强度、关联稳定性、关联时序性、分布符合性、医学和生物学合理性、剂量-反应关系。需注意\imp{混杂因素（Confounding Factor）}的控制。

\begin{table}[htbp]
\centering
\caption{环境流行病学主要研究方法}
\begin{tabularx}{\textwidth}{lXl}
\hline
\textbf{效应类型} & \textbf{研究方法} & \textbf{特点} \\
\hline
\multirow{3}{*}{急性健康效应} & 时间序列研究 & 分析暴露与效应的时间变化关系 \\
& 病例交叉研究 & 以病例自身为对照 \\
& 固定群组研究 & 追踪固定人群暴露与效应 \\
\hline
\multirow{4}{*}{慢性健康效应} & 横断面研究 & 同一时间点调查暴露与疾病 \\
& 生态学研究 & 以群体为单位分析 \\
& 病例对照研究 & 比较病例与对照暴露史 \\
& 队列研究 & 前瞻性追踪暴露与结局 \\
\hline
\end{tabularx}
\end{table}

\begin{tipbox}
\textbf{宣威女性肺癌高发研究}是环境流行病学研究的经典案例。研究发现室内燃煤空气污染（多环芳烃暴露）是主要原因，改造炉灶后肺癌发病率显著下降。启示：流行病学+干预研究的整合范式。
\end{tipbox}

\subsection{环境毒理学方法}

\begin{table}[htbp]
\centering
\caption{流行病学与毒理学方法的比较}
\begin{tabularx}{\textwidth}{lXX}
\hline
\textbf{比较维度} & \textbf{环境流行病学} & \textbf{环境毒理学} \\
\hline
研究对象 & 人群 & 实验动物/细胞 \\
优势 & 资料来自人群，结果不需外推；可用以评价干预效果；可研究对智力、心理、感觉等的影响 & 可精确控制暴露水平和条件；效应指标不受限制；可应用转基因等先进材料 \\
劣势 & 暴露难精确测量；混杂因素多；弱效应难评价；某些危害（致癌）潜伏期太长 & 种属差异（如反应停和无机砷在动物中为阴性但对人确有毒性）；通常高剂量暴露，易夸大效应 \\
结论强度 & 相关性提示 & 因果推断较强 \\
\hline
\end{tabularx}
\end{table}

\begin{tipbox}
环境流行病学和环境毒理学方法具有\imp{互补性}，开展环境与健康研究时必须将两者结合起来，相互补充、相得益彰。

\textbf{环境生物监测：}传统的环境监测主要采用化学或物理学方法测定介质中污染物的含量，而生物监测则能够迅速反映污染物是否能对生物体特别是体内的遗传物质产生影响。环境理化监测结合生物监测是今后环境监测的趋势。包括：①现场生物监测——对环境动植物或微生物进行细胞遗传学或分子毒理学的直接监测；②环境样品的生物监测——收集空气、水和固体环境样品进行毒理学测试。
\end{tipbox}

\subsection{生物标志物}

\begin{defbox}
\textbf{生物标志物（Biomarker）}：生物体内发生的与发病机制有关联的关键事件的指示物，是机体由于接触各种环境因素引起器官、细胞和亚细胞的生理、生化、免疫和遗传等任何可测定的改变。

\textbf{分类：}
\begin{enumerate}[leftmargin=*]
    \item \imp{暴露生物标志物}——包括内剂量和生物有效剂量标志
    \item \imp{效应生物标志物}——机体内可测定的生理、生化或其他方面的改变
    \item \imp{易感性生物标志物}——能指示机体接触某种特定环境因子时反应能力的标志
\end{enumerate}

\textbf{分子生物标志（Molecular Biomarkers）}：研究外来因子与生物大分子（核酸、蛋白质）相互作用引起的一切分子水平上的改变。以分子生物标志建立的分子流行病学可准确反映暴露与反应的关系，对早期预测损害、评价危险度有重大意义。

\textbf{在环境流行病学中的应用：}暴露的精确测量；揭示早期生物效应；判断宿主易感性。
\end{defbox}

% ----- 2.6 预防对策 -----
\section{环境与健康问题的预防对策}

主要原则：\imp{控制危险因素} → \imp{减少人群接触} → \imp{保护易感人群} → \imp{提升自我防护能力}

主要措施：组织措施、规划措施、技术措施、教育措施。

% ----- 2.7 标准体系 -----
\section{环境与健康标准体系}

我国环境与健康标准体系分为由\imp{环境保护部门}牵头制定的\imp{环境保护标准体系}和由\imp{卫生部门}牵头制定的\imp{环境卫生标准体系}。

\textbf{环境保护标准体系：}以保护人的健康和生存环境、防止生态环境遭受破坏为目的。按级别分三级：国家标准、地方标准（效率高于国家标准）、行业标准。按内容分五类：①环境质量标准（体系核心）；②污染物排放（控制）标准；③环境基础标准；④环境监测分析方法标准；⑤环境样品标准。按效力分：强制性环境标准和推荐性环境标准（不具有法律强制力）。

\textbf{环境卫生标准体系：}以保护人群身体健康为直接目的。包括环境卫生专业基础标准和环境卫生单项标准（室内空气、生活饮用水、公共场所、化妆品等）。

\begin{defbox}
\textbf{基准（Criterion）}：根据环境中有害物质与机体之间的剂量-反应关系，考虑敏感人群和暴露时间而确定的对健康不会产生直接或间接有害影响的相对安全剂量。

\textbf{标准（Standard）}：以保护人群健康为直接目的，对环境中有害影响的因素提出的限量要求以及实现这些要求所规定的相应措施。具有\imp{法律强制性}。
\end{defbox}

\begin{keybox}
\textbf{标准与基准的区别：}

\begin{table}[H]
\centering
\caption{基准与标准的比较}
\begin{tabularx}{\textwidth}{lXX}
\hline
\textbf{比较维度} & \textbf{基准（Criterion）} & \textbf{标准（Standard）} \\
\hline
定义 & 根据剂量-反应关系，考虑敏感人群和暴露时间，确定的对健康不产生有害影响的最大剂量或浓度 & 以保护人群健康为直接目的，对环境中有害因素提出的限量要求及实现这些要求规定的措施 \\
关系 & 标准的科学依据（核心） & 基准内容的实际体现 \\
法律效力 & 不考虑经济、社会、技术因素，无法律效力 & 考虑经济、社会、技术因素，有法律效力 \\
\hline
\end{tabularx}
\end{table}

\begin{itemize}[leftmargin=*]
    \item 基准是纯科学数据，无法律效力
    \item 标准基于基准制定，但需考虑经济、技术可行性
    \item 标准具有法律强制力，是执法依据
    \item 原则上标准限值 $\leq$ 基准值（更严格），但可有例外
\end{itemize}
\end{keybox}

\textbf{制定环境卫生标准的原则：}
\begin{enumerate}[leftmargin=*]
    \item 不引起急性或慢性中毒及潜在的远期危害
    \item 对主观感觉无不良影响
    \item 对人体健康无间接不良影响
    \item 选用最敏感指标（最敏感人群、最敏感观察指标、最灵敏方法）
    \item 经济合理和技术可行
\end{enumerate}

\begin{tipbox}
\textbf{环境卫生标准制定的方法（$\approx$基准的研究方法）：}
\begin{enumerate}[leftmargin=*]
    \item \imp{环境毒理学方法}：明确污染物的作用性质、特点及阈剂量。需考虑安全系数以调整：动物与人种属差异、人群易感性差异、高剂量短时间与低剂量长期暴露差异、小样本外推大群体差异
    \item \imp{感官机能影响的测定}：确定对眼、口腔、上呼吸道的刺激作用阈和异常颜色、气味的阈浓度。在确保安全条件下人体实验比动物实验意义更大
    \item \imp{环境流行病学研究}：验证阈浓度，分析污染与危害的相关程度，评价最高容许浓度的安全程度
    \item 其他：人类受控实验（志愿者）、人类负荷量测定、数学模型、危险度评价
\end{enumerate}
\end{tipbox}

% ----- 2.8 健康风险评价 -----
\section{健康危险度评价}

\begin{keybox}
\textbf{健康危险度评价（Health Risk Assessment, HRA）}：按一定准则，应用毒理学研究和流行病学调查等资料，系统科学地表征有害环境因素暴露对人类和生态的潜在损害作用，并对产生这种损害作用证据的强度或充分性进行评定，对危险性评估相关的不确定性进行评价。

\textbf{特点：}①健康保护观念的转变——安全是相对的，不可能将有害健康的污染物完全清除，只能逐步控制使之处于可接受水平；②把环境污染对健康的影响定量化。

\textbf{四个阶段：}
\begin{enumerate}[leftmargin=*]
    \item \imp{危害鉴定（Hazard Identification）}——定性评价，判断某物质是否对健康有害。依据主要来自流行病学（说服力最强）和毒理学
    \item \imp{暴露评价（Exposure Assessment）}——关键步骤，评估人群暴露的强度、频率和持续时间
    \item \imp{剂量-反应关系评价（Dose-Response Assessment）}——核心，量化剂量与效应关系。有阈化学物采用NOAEL法或基准剂量法推导参考剂量；无阈化学物通过数学模型外推低剂量风险
    \item \imp{风险特征分析（Risk Characterization）}——综合前三阶段，分析判断人群发生某种危害的可能性大小，并阐述不确定性
\end{enumerate}

\textbf{可接受危险度：}社会公认的可接受不良健康效应的概率。有阈化学物参考剂量对应的可接受危险度定为$10^{-6}$。

\textbf{致癌强度系数（CPF）：}实验动物或人终生暴露于剂量为每日每公斤体重1mg致癌物时的终生超额患癌危险度，为剂量-反应曲线斜率的95\%上限。
\end{keybox}

\begin{tipbox}
\textbf{致癌物评价的不确定性来源：}
\begin{enumerate}[leftmargin=*]
    \item 种属差异——不同物种间生理和代谢差异
    \item 高剂量向低剂量外推——高剂量下的反应不一定适用于低剂量
    \item 较短染毒时间向长期持续接触外推
    \item 暴露环境不同——实验暴露条件与人类实际暴露不同
    \item 毒代/毒效动力学资料不足
    \item 数据质量和可靠性问题
\end{enumerate}
\end{tipbox}

% ----- 2.9 复习题 -----
\section{复习题}

\begin{enumerate}[leftmargin=*, itemsep=8pt]
    \item \textbf{【名词解释】}生态系统；食物链；生物放大作用；生物半减期
    \item \textbf{【名词解释】}人群健康效应谱；高危人群；Hormesis效应
    \item \textbf{【名词解释】}生物标志物；环境内分泌干扰物；环境基因交互作用
    \item \textbf{【名词解释】}健康危险度评价（HRA）；可接受危险度
    \item \textbf{【名词解释】}基准与标准；环境生物监测
    \item \textbf{【名词解释】}无阈值化合物与阈值化合物（单阈值与双阈值）
    \item \textbf{【简答题】}简述人群健康效应谱的五级分布规律及其公共卫生学意义。
    \item \textbf{【简答题】}试比较环境流行病学与环境毒理学在研究环境污染健康效应中的优缺点。
    \item \textbf{【简答题】}简述健康危险度评价的基本内容和四个阶段。
    \item \textbf{【简答题】}简述制定环境卫生标准的原则及其制定方法。
    \item \textbf{【简答题】}简述环境多因素联合作用的五种类型及其定义，并说明研究联合作用的意义。
    \item \textbf{【非标准答案题】}试论生物放大作用对环境污染物健康危害评估的启示。
\end{enumerate}

\subsection*{参考答案要点}

\begin{enumerate}[leftmargin=*, itemsep=5pt]
    \item \textbf{生态系统}：在一定空间范围内，由生物群落及其环境组成，借助物质流、能量流、物种流和信息流联结的稳态系统。\textbf{食物链}：不同营养级生物逐级被吞食的链锁关系。\textbf{生物放大作用}：重金属和难降解有机物通过食物链转移到高位营养级生物体内，浓度逐级放大的现象。
    \item \textbf{人群健康效应谱}：人群暴露于污染物时产生的不良反应在性质、程度和范围上表现的征象，从弱到强分为：负荷增加→生理代偿→亚临床状态→临床疾病→死亡。呈金字塔形分布。\textbf{高危人群}：敏感人群和高暴露人群。
    \item \textbf{生物标志物}：可测量的反映暴露、效应或易感性的生物学指标。\textbf{EDCs}：具有类似激素作用，干扰内分泌功能的外源性物质。
    \item HRA：对有害环境因素作用于特定人群的有害健康效应进行综合定性、定量评价的过程，包括危害鉴定、暴露评价、剂量-反应关系评价和风险特征分析四个阶段。
    \item \textbf{基准}：基于剂量-反应关系的纯科学数据，无法律效力。\textbf{标准}：以保护健康为目的的限量要求，具有法律强制力。
    \item \textbf{无阈化合物}：任何剂量均可能有害，无安全剂量（如遗传毒性致癌物）。\textbf{阈值化合物}：仅达到或超过阈剂量才产生效应。
    \item （答题要点：五级分布的具体内容；呈金字塔分布；公共卫生学意义包括早期发现高危人群、指导标准制定、评估干预效果等）
    \item （答题要点：流行病学——优势是结果来自人群无需外推，不足是暴露难精确测量、混杂因素多；毒理学——优势是可精确控制暴露条件，不足是种属差异和高剂量外推问题）
    \item （答题要点：危害鉴定、暴露评价、剂量-反应关系评价、风险特征分析四阶段）
    \item （答题要点：不引起急慢性中毒及远期危害；对主观感觉无不良影响；对健康无间接不良影响；选用最敏感指标；经济合理技术可行）
    \item 四种类型：相加（联合=各单和）、协同（联合>各单和）、增强（无毒物增强有毒物）、拮抗（联合<各单和）。
    \item （开放性问题，从污染物通过食物链传递浓缩、不同营养级暴露风险差异、对风险评价模型的影响等方面展开。）
\end{enumerate}

\newpage

% =====================================================================
%  CHAPTER 3: 室外空气环境与健康 【重点章节】
% =====================================================================
\newpage